%% kapa: exam=1 round=27 year=2016 period=1 form=A subject=civil_law subject_ko=민법 seq=22 number=22
%% source: Q-Net 감정평가사 기출문제 (cst00309) artlSeq=5169402 file=extracted/1008/5169402/제27회 감정평가사 1차 문제/제27회 감정평가사 1교시문제(A형).pdf
%% conversion: kordoc 4.24.51 -> tools/gen1.py (auto); review=proofread-sonnet; flags=none
\begin{question}{22}
명인방법에 관한 설명으로 옳지 않은 것은? (다툼이 있으면 판례에 따름)
\begin{choices}
  \item 관습법상의 공시방법이다.
  \item 수확되지 아니한 농작물에 대해서도 인정된다.
  \item 토지의 지상물이 독립된 물건이며 현재의 소유자가 누구라는 것을 명시하여야 한다.
  \item 명인방법으로 양도담보를 공시할 수 없다.
  \item 건물 이외의 지상물을 토지 또는 원물과 분리하지 않은 채 독립된 거래객체로 하는 데 이용된다.
\end{choices}
\end{question}
\answer{4}
